\chapter{Resultados}
\label{chap:resultados}

% TODO(ensaio definitivo): substituir as secoes 4.2-4.4 pelos resultados
% das series finais (bateria nova, 21 degraus, 3 repeticoes, subida+descida).

\section{Atividades do Projeto}
\label{sec:atividades}

As atividades desenvolvidas compreenderam: (i) estudo do sistema de
propulsao e do ESC; (ii) selecao e validacao em bancada do sensor de
corrente; (iii) construcao da bancada de tracao; (iv) implementacao da
aquisicao automatica (Arduino Due + Raspberry Pi); (v) ensaios preliminares
de caracterizacao; e (vi) analise dos dados com tratamento de incertezas.

\section{Validação do Sensor de Corrente}
\label{sec:resultsensor}

Em bancada, com fonte de 0 a 6 A e multimetro de referencia, o sistema de
aquisicao com ACS758 apresentou relacao linear com a corrente aplicada em
toda a faixa, com regressao $y = 1{,}0742\,x$ entre os valores medidos pelo
multimetro ($x$) e pelo sistema ($y$). Os desvios sistematicos de $\pm$7\%
sao compensaveis por fator de correcao em software, validando o sensor para
monitoramento.

\section{Curva Característica Torque $\times$ Corrente (Ensaios Preliminares)}
\label{sec:curva}

O ensaio preliminar de tração (escada de 97$^\circ$ a 115$^\circ$,
superficie com EVA) caracterizou a curva na faixa de 3,6 a 18,2 A. Na
janela estacionaria (7--10 s de cada degrau), o ajuste por minimos
quadrados resultou em
\begin{equation}
T \;[\text{N$\cdot$m}] = 0{,}185 \; I_{ef} \;[\text{A}] + 0{,}48,
\qquad r = 0{,}92,
\end{equation}
com torque crescente de aproximadamente 1,3 N$\cdot$m (3,6 A) a
3,5 N$\cdot$m (18,2 A). A forca de tracao nao apresentou saturacao por
aderencia ate o ultimo degrau, indicando que a escada pode ser estendida
para caracterizacao da faixa de correntes maiores.

\subsection{Validação Cruzada com o Modelo do Powertrain}

Comparando degrau a degrau com o modelo teorico
($K_t$ do motor de 2100 KV, reducao $i \approx 11{,}5$, $\eta$ estimado),
os desvios ficaram dentro de $\pm$5\% nos degraus de 104$^\circ$ a
115$^\circ$, incluindo concordancia exata (0,0\%) no degrau maximo,
confirmando a consistencia da cadeia de medicao e do modelo fisico.
Degraus proximos ao neutro foram excluidos da comparacao pela incerteza do
\emph{duty cycle} efetivo na regiao nao-linear do ESC.

\section{Testes}
\label{sec:testes}

Os ensaios preliminares revelaram efeitos sistematicos tratados no
protocolo: (i) corte de baixa tensao do ESC por bateria degradada
(sesssao descartada e repetida com bateria saudavel); (ii) perda de aderencia
na descida (resolvida com massa adicional e limpeza dos pneus); e
(iii) dessincronizacao entre comando e registro (eliminada com espera de
finalizacao do comando em cada degrau).

% TODO(ensaio definitivo):
% \section{Analise dos Resultados} -> regressao final, resíduos,
% repetitividade (t-Student), incerteza expandida, linearidade, histerese.
% Incluir figuras: curva T x I com banda de incerteza; residuos.

\section{Resumo do Capítulo}

O sensor de corrente foi validado em bancada e a metodologia de medicao
indireta de torque produziu uma curva caracteristica coerente com o modelo
do powertrain na faixa preliminar. Os ensaios definitivos ampliarao a faixa
de corrente e completarao a caracterizacao metrologica.

\clearpage
